\documentclass{article}

\title{A Short Demonstration of \LaTeX.js}
\author{Liton}
\date{2026}

\begin{document}

\maketitle

\begin{abstract}
This page is rendered live in your browser by \LaTeX.js, a \LaTeX{} to HTML5
translator written in JavaScript. It shows a small but complete document:
sections, lists, a table, mathematics, and typography.
\end{abstract}

\section{Introduction}

\LaTeX.js parses \LaTeX{} source and produces HTML together with a stylesheet
that mimics the default \LaTeX{} output. Unlike a math-only renderer such as
KaTeX or MathJax, it understands the document \emph{structure} as well:
sections, lists, tables, and floating figures.

The idea is simple. You write ordinary \LaTeX{}, and the result is a normal
web page --- no \TeX{} installation, no build step on the server.

\subsection{What this document covers}

The constructs below are the ones you would normally reach for in a short
report:

\begin{itemize}
  \item sectioning: sections and subsections
  \item lists: both unordered and ordered
  \item tables
  \item mathematics: inline and displayed
  \item typography: \textbf{bold}, \emph{emphasis}, \texttt{typewriter}
\end{itemize}

\section{Lists}

An ordered list behaves as expected:

\begin{enumerate}
  \item Parse the \LaTeX{} source.
  \item Build the document tree.
  \item Emit HTML and apply the stylesheet.
\end{enumerate}

\section{Tables}

\begin{center}
\begin{tabular}{|l|c|r|}
\hline
Project & Kind & Bundle size \\
\hline
\LaTeX.js & full document & 593 KB \\
KaTeX & math only & 280 KB \\
MathJax & math only & 1.1 MB \\
\hline
\end{tabular}
\end{center}

The table above lists three projects that translate \LaTeX{} for the web,
together with what each of them is able to handle.

\section{Mathematics}

Inline mathematics sits in the flow of a sentence: for instance
$\sum_{i=1}^{n} i = \frac{n(n+1)}{2}$ is the familiar closed form.

A displayed equation is set on its own line:

$$
  \int_{0}^{1} x^{2} \, dx = \frac{1}{3}
$$

Subscripts and superscripts compose naturally, as in
$a^{2} + b^{2} = c^{2}$, and Greek letters such as $\alpha$, $\beta$ and
$\gamma$ are available.

\section{Typography}

The Computer Modern fonts are shipped with the renderer, so the output really
does look like a \LaTeX{} document rather than a web page pretending to be
one. Ligatures, spacing and the shape of the letters all come from the same
typefaces \TeX{} uses.

\section{Conclusion}

Everything on this page was produced from plain \LaTeX{} source, inside the
browser, with no server-side processing at all.

\end{document}
